\chapter{Changement de \textit{hostname} et renommage de l'utilisateur principal}

\section*{Introduction}

Cette partie s'adresse aux utilisateurs d'Ubuntu Server, ayant créé une image d'Ubuntu dans le but de la copier sur une nouvelle Raspberry Pi.  
Nous verrons comment modifier de façon durable le hostname, car il est nécessaire d'en avoir des différents sur chaque robot pour utiliser les programmes du workspace \textit{mecanum}.  
Nous verrons également comment renommer proprement l'utilisateur principal pour correspondre au nouveau hostname.



\section{Création d'un utilisateur temporaire}

Avant de modifier le hostname principal ou de renommer un utilisateur, il est recommandé de créer un utilisateur temporaire disposant des droits \textit{sudo} pour ne pas perdre l'accès administrateur en cas de problème.

\begin{lstlisting}[style=commandstyle]
sudo adduser tempadmin
sudo usermod -aG sudo tempadmin
\end{lstlisting}

Déconnectez-vous, puis reconnectez-vous en tant que \textit{tempadmin} pour poursuivre les modifications.  
Pour cela, vous pouvez soit redémarrer la Raspberry Pi, soit utiliser les commandes suivantes :

\begin{lstlisting}[style=commandstyle]
exit
ssh tempadmin@ip_de_votre_machine
\end{lstlisting}



\section{Arrêter tous les processus de l'utilisateur à renommer}

Pour éviter des conflits lors du renommage, arrêtez tous les processus de l'utilisateur à modifier (ici \textit{dartagnan}) :

\begin{lstlisting}[style=commandstyle]
sudo pkill -u dartagnan
\end{lstlisting}


\section{Renommer l'utilisateur principal}

\subsection*{1. Renommer le login}

\begin{lstlisting}[style=commandstyle]
sudo usermod -l mercedes dartagnan
\end{lstlisting}

\subsection*{2. Renommer le dossier personnel}

\begin{lstlisting}[style=commandstyle]
sudo usermod -d /home/mercedes -m mercedes
\end{lstlisting}

\subsection*{3. Renommer le groupe principal (optionnel)}

\begin{lstlisting}[style=commandstyle]
sudo groupmod -n mercedes dartagnan
\end{lstlisting}

\subsection*{4. Assurer les bonnes permissions sur le dossier personnel}

\begin{lstlisting}[style=commandstyle]
sudo chown -R mercedes:mercedes /home/mercedes
\end{lstlisting}


\section{Modifier la configuration de \textit{cloud-init}}

Commencez par éditer le fichier de configuration principal de \textit{cloud-init} :

\begin{lstlisting}[style=commandstyle]
sudo nano /etc/cloud/cloud.cfg
\end{lstlisting}

Localisez la ligne contenant \textit{preserve\_hostname: false} et remplacez-la par :

\begin{lstlisting}[style=commandstyle]
preserve\_hostname: true
\end{lstlisting}

Cela indique à \textit{cloud-init} de ne pas modifier le hostname au démarrage. Sauvegardez et quittez l'éditeur.


\section{Modifier les fichiers \textit{/etc/hostname} et \textit{/etc/hosts}}

Modifiez le fichier \textit{/etc/hostname} pour indiquer le nouveau nom souhaité (par exemple \textit{mercedes-desktop}) :

\begin{lstlisting}[style=commandstyle]
sudo nano /etc/hostname
\end{lstlisting}

Changez le contenu en remplaçant l'ancien hostname par le nouveau, puis sauvegardez.

Modifiez aussi le fichier \textit{/etc/hosts} pour mettre à jour la correspondance avec le nouveau hostname :

\begin{lstlisting}[style=commandstyle]
sudo nano /etc/hosts
\end{lstlisting}

Recherchez la ligne contenant l'ancien hostname, par exemple :

\begin{lstlisting}[style=commandstyle]
127.0.1.1    mercedes
\end{lstlisting}

et remplacez-la par :

\begin{lstlisting}[style=commandstyle]
127.0.1.1    mercedes-desktop
\end{lstlisting}

Sauvegardez et fermez l'éditeur.


\section{Vérification optionnelle de la configuration Netplan}

Bien que rarement problématique, il est utile de vérifier si Netplan ne force pas un hostname au démarrage.

Listez les fichiers de configuration Netplan :

\begin{lstlisting}[style=commandstyle]
ls /etc/netplan/
\end{lstlisting}

Ouvrez le fichier YAML principal (le nom peut varier, ici \textit{50-cloud-init.yaml}) :

\begin{lstlisting}[style=commandstyle]
sudo nano /etc/netplan/50-cloud-init.yaml
\end{lstlisting}

Si vous trouvez une ligne \textit{set-hostname} ou \textit{hostname}, commentez-la ou supprimez-la.



\section{Redémarrage}

Pour appliquer toutes les modifications, redémarrez la machine :

\begin{lstlisting}[style=commandstyle]
sudo reboot
\end{lstlisting}

Après redémarrage, votre invite de commande devrait refléter le nouveau hostname et utilisateur, par exemple :

\begin{lstlisting}[style=commandstyle]
mercedes@mercedes-desktop:~$
\end{lstlisting}

